% presentation/slides/05_vraisemblance.tex
\begin{frame}{Le Défi de la Log-Vraisemblance Marginale}
    Puisque les variables $Z_i$ sont manquantes, on doit intégrer (marginaliser) la densité conjointe :
    
    $$l_n(\theta) = \sum_{i=1}^n \log \left( \int_{\mathcal{Z}} f_Y(Y_i | Z=z, \theta) f_Z(z) d\mu(z) \right)$$
    
    \vspace{0.2cm}
    \textbf{Une structure non linéaire complexe (Log-Sum-Exp) :}
    \begin{itemize}
        \item $l_n(\theta)$ \alert{\textbf{n'est plus concave}}, invalidant une optimisation exacte globale directe.
        \item Présence de multiples \textbf{maximums locaux} (pièges pour les algorithmes).
        \item \textbf{Défaut d'identifiabilité :} Le problème est invariant à une translation globale près. On doit introduire une \textbf{distance quotient d'orbite} :
        $$d_{\text{orbite}}(\theta, \theta^*) = \min_{z \in \mathcal{Z}} \| \theta - A_z \theta^* \|$$
    \end{itemize}
\end{frame}